\small
\begin{xltabular}{\textwidth}{@{}l Y l c c@{}}
  \caption{Relación de historias de usuario del proyecto, agrupadas por fase.}
  \label{tab:backlog} \\
  \toprule
  \textbf{ID} & \textbf{Historia} & \textbf{Tipo} & \textbf{Esfuerzo} & \textbf{Prioridad} \\
  \midrule
  \endfirsthead
  \multicolumn{5}{@{}l}{\small\textit{(continuación de la página anterior)}} \\
  \toprule
  \textbf{ID} & \textbf{Historia} & \textbf{Tipo} & \textbf{Esfuerzo} & \textbf{Prioridad} \\
  \midrule
  \endhead
  \hline
  \multicolumn{5}{r@{}}{\small\textit{(continúa en la página siguiente)}} \\
  \endfoot
  \bottomrule
  \endlastfoot
  \multicolumn{5}{@{}l}{\textbf{Fase 0 — Cimientos: datos, scraping y fragmentación}} \\[2pt]
  IT-01 & Modelos de datos y validadores & Código & S & MUST \\
  IT-02 & Limpiador de texto (TextCleaner) & Código & S & MUST \\
  IT-03 & Spider: listado de grados de la EPSJ & Código & S & MUST \\
  IT-04 & Spider: portada de grado y detección de doble grado & Código & S & MUST \\
  IT-05 & Spider: tabla híbrida de asignaturas & Código & M & MUST \\
  IT-06 & Spider: extracción de guía docente & Código & M & MUST \\
  IT-07 & Spider: salidas profesionales & Código & S & MUST \\
  IT-08 & Chunking semántico base & Código & M & MUST \\
  IT-09 & Chunking: fusión de chunks pequeños y filtrado & Código & S & MUST \\
  IT-10 & Script de verificación del dataset & Código & S & MUST \\
  IT-11 & ADR-0001: estrategia de chunking & ADR & S & MUST \\
  IT-12 & ADR-0002: elección de Scrapy frente a alternativas & ADR & S & MUST \\
  IT-13 & README + empaquetado del proyecto & Gestión & XS & MUST \\
  IT-14 & Memoria Cap.1 y Cap.3: introducción, objetivos y planteamiento & Memoria & M & MUST \\
  IT-15 & Memoria Cap.4: alternativas de scraping & Memoria & S & SHOULD \\
  IT-16 & Memoria Cap.4: alternativas de chunking & Memoria & S & SHOULD \\
  IT-17 & Memoria Cap.4: diseño de la arquitectura del pipeline & Memoria & M & MUST \\
  IT-63 & Memoria Cap.4: implementación del spider de extracción & Memoria & M & MUST \\
  IT-64 & Memoria Cap.4: implementación del chunker y deduplicación & Memoria & S & MUST \\
  IT-19 & Memoria Cap.4: verificación y reproducibilidad de la Fase 0 & Memoria & S & MUST \\
  IT-20 & Estudio de alternativas y viabilidad & Memoria & M & SHOULD \\
  IT-21 & Hipótesis y restricciones & Memoria & XS & SHOULD \\
  IT-22 & Requisitos de usuario & Memoria & S & SHOULD \\
  IT-23 & Diagramas de secuencia: scraping y chunking & Memoria & M & SHOULD \\
  IT-24 & Modelo de datos del sistema & Memoria & S & SHOULD \\
  IT-25 & Diagrama de Gantt y captura del tablero Kanban & Memoria & XS & SHOULD \\
  IT-26 & Glosario de definiciones y abreviaturas & Memoria & S & SHOULD \\
  IT-55 & Pipeline CI/CD: validación semanal del dataset & Gestión & S & MUST \\
  IT-56 & Deduplicación de guías compartidas entre titulaciones & Código & M & MUST \\
  IT-57 & Sincronizar estructura de memoria desde Overleaf & Gestión & XS & SHOULD \\
  IT-60 & Quitar del CI la verificación de un dataset no versionado & Gestión & XS & MUST \\
  IT-61 & Correcciones de la primera revisión del tutor sobre la memoria & Memoria & M & MUST \\
  IT-62 & Segunda ronda de revisión de la memoria & Memoria & M & SHOULD \\
  IT-65 & DQA-0001: anomalías de la fuente de datos EPSJ & DQA & S & MUST \\
  \addlinespace
  \multicolumn{5}{@{}l}{\textbf{Fase 1 — Incrustaciones y base de datos vectorial}} \\[2pt]
  IT-27 & Diseño del conjunto de evaluación (preguntas etiquetadas) & Experimento & M & MUST \\
  IT-28 & Experimento comparativo de modelos de embeddings & Experimento & L & MUST \\
  IT-29 & ADR-0003: elección de modelo de embeddings & ADR & S & MUST \\
  IT-30 & Pipeline de indexación & Código & M & MUST \\
  IT-31 & Experimento comparativo de vector DB & Experimento & M & MUST \\
  IT-32 & ADR-0004: elección de vector DB & ADR & S & MUST \\
  IT-33 & Memoria Cap.4: desarrollo de la Fase 1 & Memoria & M & MUST \\
  IT-59 & Requisitos de sistema (especificación técnica) & Memoria & S & SHOULD \\
  IT-66 & Revisión macro/micro de la memoria al cerrar la Fase 0 & Memoria & S & SHOULD \\
  IT-67 & Spider: detectar respuestas que no son HTML (guías en PDF) & Código & M & MUST \\
  IT-68 & Memoria Cap.4: canalización de indexación & Memoria & S & MUST \\
  IT-69 & Memoria Cap.5: diseño del conjunto de evaluación & Memoria & S & MUST \\
  \addlinespace
  \multicolumn{5}{@{}l}{\textbf{Fase 2 — Pipeline RAG y evaluación}} \\[2pt]
  IT-34 & Diseño del prompt de generación & Código & S & MUST \\
  IT-35 & Elección y configuración del LLM de generación & Experimento & M & MUST \\
  IT-36 & ADR-0005: elección del LLM de generación & ADR & S & MUST \\
  IT-37 & Implementación del pipeline RAG completo & Código & L & MUST \\
  IT-38 & Evaluación cuantitativa del sistema RAG & Experimento & L & MUST \\
  IT-39 & Amenazas a la validez del experimento de evaluación & Experimento & S & MUST \\
  IT-40 & Memoria Cap.4/5: desarrollo de la Fase 2 y evaluación experimental & Memoria & L & MUST \\
  IT-41 & Diagrama de secuencia del flujo RAG completo & Memoria & S & SHOULD \\
  IT-58 & Memoria Cap.5: caracterización del corpus de la Fase 0 y amenazas a su validez & Memoria & S & MUST \\
  \addlinespace
  \multicolumn{5}{@{}l}{\textbf{Fase 3 — Aplicación web de chat}} \\[2pt]
  IT-42 & Diseño de la arquitectura de la app web & Código & S & SHOULD \\
  IT-43 & Storyboards de la interfaz de chat & Memoria & S & SHOULD \\
  IT-44 & Backend: endpoint de chat conectado al RAG & Código & M & SHOULD \\
  IT-45 & Frontend de chat & Código & M & SHOULD \\
  IT-46 & Memoria Cap.4: desarrollo de la Fase 3 (app web) & Memoria & S & SHOULD \\
  IT-47 & Anexo: documentación técnica de la API & Memoria & S & COULD \\
  \addlinespace
  \multicolumn{5}{@{}l}{\textbf{Fase 4 — Validación y estudio de ablación}} \\[2pt]
  IT-48 & Validación con usuarios reales & Experimento & M & COULD \\
  IT-49 & Estudio de ablation & Experimento & M & SHOULD \\
  IT-50 & Memoria: retroalimentación de usuarios y conclusiones & Memoria & M & SHOULD \\
  IT-51 & Anexo: resultados completos de la validación con usuarios & Memoria & S & COULD \\
  \addlinespace
  \multicolumn{5}{@{}l}{\textbf{Fase 5 — Cierre y defensa}} \\[2pt]
  IT-52 & Revisión formal completa de la memoria & Memoria & M & MUST \\
  IT-53 & Material de apoyo para la defensa + ensayo oral & Gestión & S & SHOULD \\
  IT-54 & Anexo: instalación y configuración del sistema & Memoria & S & COULD \\
\end{xltabular}
\normalsize
